% GENERATED FILE. Edit canonical lessons, then run npm run generate:books.
% canonical-source-hash: fnv1a64:26831423004dd34f
% canonical-lessons: GE-C20-entschuldigung, GE-R20-danke-bitte-entschuldigung

\chapter{Sorry}
\label{ch:ge-sorry}
\hlchaptermodality{33}

\begin{chapteropening}
\textbf{By the end of this chapter:} \emph{I can apologise or excuse myself in German with Entschuldigung, and reach for es tut mir leid when I mean it more.}

Say Schuld, then Entschuldigung as the 'un-guilting' built on it, and switch to es tut mir leid for real regret rather than a bump in the hallway. The chapter's only lesson, so the payoff covers everything it introduces.
\end{chapteropening}

\section[Entschuldigung]{Entschuldigung — “sorry” / “excuse me”}

\label{lesson:GE-C20-entschuldigung}



\begin{quote}
You already know \textbf{bitte} does triple duty for \textquotedblleft{}please,\textquotedblright{} \textquotedblleft{}you're welcome,\textquotedblright{} and \textquotedblleft{}here you go.\textquotedblright{} German's word for \textbf{sorry} is just as hard-working: \textbf{Entschuldigung} covers both \textquotedblleft{}\textbf{excuse me}\textquotedblright{} and \textquotedblleft{}\textbf{I'm sorry}.\textquotedblright{}
\end{quote}

\begin{cousinweb}
\begin{itemize}
\raggedright
  \item \textbf{ent-} = a \textquotedblleft{}removal / undoing\textquotedblright{} prefix — the same shape as in \textbf{ent}decken (\textquotedblleft{}un-cover\textquotedblright{} $\to$ discover) or \textbf{ent}fernen (\textquotedblleft{}un-distance\textquotedblright{} $\to$ remove).
  \item \textbf{Schuld} = \textquotedblleft{}\textbf{guilt, fault, blame}\textquotedblright{} (also, separately, \textquotedblleft{}debt\textquotedblright{} — a \emph{Schuld} can be money owed or wrong done).
  \item \textbf{-igung} = a noun-forming suffix, roughly \textquotedblleft{}-ing/-ation.\textquotedblright{}
\end{itemize}

So \textbf{Entschuldigung} is literally \textquotedblleft{}\textbf{an un-guilting}\textquotedblright{} — the act of removing blame. Saying it hands the guilt back: \emph{I release you from any fault} (when getting someone's attention or brushing past) or \emph{I ask you to release me from mine} (when apologizing).

\textbf{Schuld} has a striking cousin: it traces to the same ancient Germanic root as English \textbf{shall} and \textbf{should} — a verb that originally meant \textbf{\textquotedblleft{}to owe.\textquotedblright{}} \emph{Schuld} (what is owed, guilt) and \emph{shall} (what one is obliged to do) are, at root, the same idea of an obligation outstanding.
\end{cousinweb}

\begin{culture}
For \textbf{deeper regret} — real sorrow, not a bump in the hallway — German often reaches for \textbf{es tut mir leid}, literally \textquotedblleft{}\textbf{it does to-me sorrow}\textquotedblright{}:

\begin{itemize}
\raggedright
  \item \textbf{es} = \textquotedblleft{}it,\textquotedblright{} \textbf{tut} = \textquotedblleft{}does\textquotedblright{} (from \textbf{tun}, \textquotedblleft{}to do\textquotedblright{}), \textbf{mir} = \textquotedblleft{}to me,\textquotedblright{} \textbf{Leid} = \textquotedblleft{}sorrow, suffering\textquotedblright{} (a cousin of English \textbf{loath}, both from a root meaning \textquotedblleft{}hateful, painful\textquotedblright{}).
\end{itemize}

Use \textbf{Entschuldigung} for the everyday sorry/excuse-me (stepping on a foot, interrupting), and \textbf{es tut mir leid} when you mean it more — condolences, a real regret. The two overlap; native speakers move between them by feel.
\end{culture}

\subsection*{Your turn}
Say these aloud:

\begin{itemize}
\raggedright
  \item \textquotedblleft{}Schuld\textquotedblright{} — guilt/fault — then \textquotedblleft{}Entschuldigung,\textquotedblright{} sorry / excuse me
  \item the deeper phrase — \textquotedblleft{}es tut mir leid,\textquotedblright{} it pains me / I'm so sorry
\end{itemize}

\begin{tcolorbox}[breakable,colback=teal!4,colframe=teal!35!black,title={Before you move on}]
What does \textbf{Entschuldigung} literally mean, piece by piece? (\textbf{\textquotedblleft{}Un-guilting\textquotedblright{}} — \emph{ent-} \textquotedblleft{}away\textquotedblright{} + \emph{Schuld} \textquotedblleft{}guilt/fault.\textquotedblright{}) What English word shares Schuld's ancient root? (\textbf{Shall/should} — both from a root meaning \textquotedblleft{}to owe.\textquotedblright{}) What's the deeper phrase for real regret? (\textbf{Es tut mir leid}, \textquotedblleft{}it does to me sorrow.\textquotedblright{})
\end{tcolorbox}
\section[Practice]{danke, bitte, Entschuldigung — and the sizes each of them comes in}

\label{lesson:GE-R20-danke-bitte-entschuldigung}



\begin{quote}
Thank someone. Then say the word that answers it. Both are two syllables ending in the same unstressed \textbf{-e} — say them back to back and let that ending stay weak.
\end{quote}

\subsection*{Your turn: three sizes, smallest first}
Say these aloud:

\begin{itemize}
\raggedright
  \item the plain one — \emph{danke}
  \item the one that upgrades it with an adjective — \emph{danke schön}
  \item the warmest everyday one — \emph{vielen Dank}
\end{itemize}

Notice what changed in the third. \emph{Danke} is a verb form — \emph{I thank} — but \emph{vielen Dank} is \textbf{the noun}, capitalised, with \emph{viel} agreeing in front of it. That is why \emph{viel} has an ending there and \emph{schön} does not: one is describing a noun, the other is not.

\subsection*{Your turn: the word that answers everything}
\textbf{bitte} is four words in one coat. Say it four times, meaning something different each time:

Say these aloud:

\begin{itemize}
\raggedright
  \item \emph{please}, asking for something
  \item \emph{you're welcome}, answering a thank-you
  \item \emph{here you go}, handing something over
  \item \emph{pardon?}, when you did not catch it
  \item and the two shortest answers there are — \emph{ja}, \emph{nein}
\end{itemize}

\begin{culture}
\textbf{Entschuldigung} is built out of \textbf{Schuld}, \emph{guilt} or \emph{fault}, with a prefix that undoes it: an \emph{un-guilting}. The word is a request that the blame be removed, which is why it works both for bumping into someone and for getting their attention in the first place.

That is worth holding against English, where \emph{sorry} and \emph{excuse me} split the job in two. German runs one word across both, and the size of the apology is carried by your voice rather than by a different word.
\end{culture}

\begin{tcolorbox}[breakable,colback=teal!4,colframe=teal!35!black,title={Before you move on}]
Why is it \emph{vielen Dank} and not \emph{viel Dank}? (\textbf{Dank is a noun, so viel takes an ending.}) What is inside \emph{Entschuldigung}? (\textbf{Schuld}, guilt — the word asks for it to be undone.) Name two things \emph{bitte} can mean.
\end{tcolorbox}
